% !TeX root = ../main.tex
\chapter{Tạo bảng, dữ liệu mẫu và truy vấn dữ liệu (10đ)}\label{chap:schema}
\section{Các câu lệnh tạo bảng và ràng buộc kiểm tra (3.5đ)}\label{sec:schema}
\subsection{Tổ chức lược đồ}
Hệ thống quản lý đồng thời thời gian sử dụng máy, ví hội viên, bán dịch vụ và tồn kho. Một sự kiện kinh tế được lưu qua chứng từ: phiên và biên lai cho tiền máy; đơn cùng dòng chi tiết cho dịch vụ; phiếu cùng dòng chi tiết cho nhập, xuất hoặc kiểm kê. Cách tổ chức này giúp truy nguyên thay đổi số dư và tồn thay vì chỉ lưu các giá trị tổng hiện tại.

MySQL 8.4.13 sử dụng InnoDB để hỗ trợ khóa ngoại và giao dịch. Các mã định danh dùng so sánh nhị phân \code{ascii_bin}; dữ liệu văn bản dùng \code{utf8mb4_0900_as_cs}. Tiền VND được lưu theo số nguyên; giá trị dẫn xuất ở dòng chi tiết được tính từ số lượng và đơn giá.
\begin{table}[H]\centering\small\caption{Năm phân hệ trong lược đồ 23 bảng}
\begin{tabularx}{\linewidth}{@{}p{34mm}Yc@{}}\toprule
\textbf{Phân hệ}&\textbf{Các bảng}&\textbf{Số bảng}\\\midrule
Danh tính & NHOM\_NGUOI\_DUNG, TAI\_KHOAN, SO\_DIEN\_THOAI, NGUOI\_CHOI, NHAN\_VIEN, DIA\_CHI\_NV & 6\\
Máy trạm & KHU\_VUC, NHOM\_MAY, BANG\_GIA, MAY\_TRAM, PHIEN\_SU\_DUNG, BIEN\_LAI\_TIEN\_MAY & 6\\
Dịch vụ & NHOM\_DICH\_VU, DON\_VI\_TINH, MAY\_IN, DICH\_VU, DON\_GOI\_DICH\_VU, CHI\_TIET\_GOI\_DICH\_VU & 6\\
Kho & PHIEU\_KHO, CHI\_TIET\_PHIEU\_KHO & 2\\
Nhân sự và ví & CA\_TRUC, PHAN\_CONG\_TRUC, GIAO\_DICH\_VI & 3\\\bottomrule
\end{tabularx}\end{table}
\subsection{Câu lệnh tạo bảng}
Các tệp sau là nguồn DDL được dùng khi khởi tạo database mới. Khóa chính định danh từng thực thể; UNIQUE ngăn trùng mã nghiệp vụ; CHECK giới hạn trạng thái, miền tiền, số lượng và thứ tự thời gian. Cột sinh kết hợp UNIQUE được dùng cho các trường hợp cần đảm bảo duy nhất có điều kiện, chẳng hạn một máy chỉ có một phiên đang hoạt động.
\sqlinput{sql/01_identity.sql}{DDL phân hệ danh tính}
\sqlinput{sql/02_workstations.sql}{DDL phân hệ máy trạm}
\sqlinput{sql/03_services.sql}{DDL phân hệ dịch vụ}
\sqlinput{sql/04_inventory.sql}{DDL phân hệ kho}
\sqlinput{sql/05_staff_wallet.sql}{DDL phân hệ nhân sự và ví}
\subsection{Khóa ngoại và chính sách toàn vẹn}
Khóa ngoại bảo đảm các dòng chi tiết và chứng từ tham chiếu đối tượng tồn tại. Các tham chiếu lịch sử dùng RESTRICT khi việc xóa đối tượng cha có thể làm mất khả năng đối soát. CASCADE chỉ được sử dụng theo quan hệ đã định nghĩa trong DDL; không được hiểu là quyền tùy ý xóa chứng từ đã chốt.
\sqlinput{sql/constraints.sql}{Các khóa ngoại được áp dụng}
\snapshot{img_01_dbms_tables}{Kết quả DBMS minh họa câu lệnh tạo bảng thực thi thành công (23 bảng)}{fig:img_01_dbms_tables}
\snapshot{img_02_dbms_constraints}{Kết quả DBMS minh họa các ràng buộc khóa ngoại và CHECK được áp dụng}{fig:img_02_dbms_constraints}
\section{Câu lệnh tạo chỉ mục (1.5đ)}\label{sec:index}
\subsection{Mô tả lý do cần đánh chỉ mục}
Ở màn hình gọi món tại máy hoặc bán tại quầy, người dùng lọc dịch vụ đang kinh doanh theo nhóm và khoảng giá, sau đó xem giá từ cao xuống thấp. Yêu cầu này xuất phát từ chức năng xem thực đơn phục vụ khách và thu ngân. Truy vấn Q-P12-02 vì vậy là cơ sở lựa chọn chỉ mục nhiều cột, không đánh chỉ mục tách rời khỏi nhu cầu sử dụng thực tế.
\sqlinput{sql/indexes.sql}{Câu lệnh tạo chỉ mục: ix\_dich\_vu\_filter}
\subsection{Minh họa hiệu quả truy vấn trước và sau khi đánh chỉ mục}
Phép đo dùng database riêng với 20.000 dòng dịch vụ, 10 nhóm và hai đơn vị tính. Cùng một câu SELECT được chạy 20 lượt trước và sau khi thêm \code{ix_dich_vu_filter}. Cả hai giai đoạn đều lấy hết tập kết quả. Số đo phụ thuộc máy chạy, cache và tải tại thời điểm đo; đây là bằng chứng cho tập dữ liệu cụ thể, không phải cam kết hiệu năng trong mọi môi trường.
\begin{table}[H]\centering\caption{So sánh hiệu năng truy vấn trước và sau khi đánh chỉ mục}\begin{tabularx}{\linewidth}{@{}Yrr@{}}\toprule\textbf{Chỉ tiêu}&\textbf{Trước}&\textbf{Sau}\\\midrule
Thời gian trung bình 20 lượt (ms)&267.313&68.820\\
Số dòng kết quả&949&949\\
Có bước Sort&Có&Có\\\bottomrule\end{tabularx}\end{table}
Hiệu quả tăng tốc đạt \textbf{3.88 lần} (từ 267.3ms xuống 68.8ms). Kế hoạch thực thi dưới đây minh họa chi tiết đường dẫn truy xuất dữ liệu từ Table Scan sang Index Range Scan:
\sqlinput{sql/plan_before.txt}{Kế hoạch thực thi EXPLAIN ANALYZE trước khi tạo chỉ mục}
\sqlinput{sql/plan_after.txt}{Kế hoạch thực thi EXPLAIN ANALYZE sau khi tạo chỉ mục}
\snapshot{index_before}{Kết quả DBMS hiển thị EXPLAIN ANALYZE trước khi tạo chỉ mục}{fig:index_before}
\snapshot{index_after}{Kết quả DBMS hiển thị EXPLAIN ANALYZE sau khi tạo chỉ mục}{fig:index_after}
\section{Truy vấn dữ liệu (5đ)}\label{sec:queries}
Yêu cầu: 5 câu truy vấn độc lập trên các bảng dữ liệu khác nhau, kèm dữ liệu mẫu và kết quả thực thi trên DBMS. Hai câu đầu có SELECT, FROM, WHERE, ORDER BY; hai câu tiếp theo có GROUP BY và AGGREGATE; câu cuối cùng lồng 2 cấp và AGGREGATE. Các bảng dữ liệu mẫu chi tiết được đối chiếu thêm tại Phụ lục \ref{app:data}.
\subsection{Câu truy vấn 1: Danh sách máy theo khu vực (Q-P12-01)}
Mô tả: Dùng cho màn hình sơ đồ máy của thu ngân để kiểm tra các máy có thể sử dụng tại KV01, hiển thị tên máy, nhóm máy và đơn giá giờ; sắp xếp theo vị trí hàng, cột.
\sqlinput{sql/Q-P12-01.sql}{Câu lệnh SELECT - Truy vấn 1 (Q-P12-01)}
\begin{table}[H]
\centering\small
\caption{Bảng dữ liệu mẫu và kết quả thực thi Truy vấn 1}
\begin{tabularx}{\linewidth}{@{}XXXXX@{}}\toprule
\textbf{Mã máy} & \textbf{Tên máy} & \textbf{Trạng thái} & \textbf{Nhóm máy} & \textbf{Giá giờ}\\\midrule
M01 & Máy 01 & AVAILABLE & Máy VIP & 15000\\
M02 & Máy 02 & AVAILABLE & Máy VIP & 15000\\
M04 & Máy 04 & AVAILABLE & Máy VIP & 15000\\
\bottomrule\end{tabularx}\end{table}
\snapshot{img_04_q01}{Kết quả màn hình hiển thị từ DBMS minh họa thực thi Truy vấn 1}{fig:img_04_q01}
\subsection{Câu truy vấn 2: Lọc thực đơn theo nhóm và giá (Q-P12-02)}
Mô tả: Dùng cho màn hình gọi món của khách hoặc bán tại quầy của thu ngân để tìm các món thuộc nhóm NDV01 có giá từ 10.000 đến 50.000đ; sắp theo giá giảm dần rồi theo mã.
\sqlinput{sql/Q-P12-02.sql}{Câu lệnh SELECT - Truy vấn 2 (Q-P12-02)}
\begin{table}[H]
\centering\small
\caption{Bảng dữ liệu mẫu và kết quả thực thi Truy vấn 2}
\begin{tabularx}{\linewidth}{@{}XXXXX@{}}\toprule
\textbf{Mã DV} & \textbf{Tên dịch vụ} & \textbf{Đơn vị} & \textbf{Giá bán} & \textbf{Tồn}\\\midrule
DV03 & Cơm chiên trứng & Phần & 25000 & 30\\
DV01 & Mì trứng & Phần & 15000 & 50\\
\bottomrule\end{tabularx}\end{table}
\snapshot{img_05_q02}{Kết quả màn hình hiển thị từ DBMS minh họa thực thi Truy vấn 2}{fig:img_05_q02}
\subsection{Câu truy vấn 3: Tổng hợp đơn đã thanh toán (Q-P12-03)}
Mô tả: Dùng cho màn hình báo cáo ca của quản lý để thống kê số đơn và doanh số theo ngày đặt và kênh đặt; có nhóm theo ngày, kênh và tính tổng tiền, đếm số đơn.
\sqlinput{sql/Q-P12-03.sql}{Câu lệnh SELECT - Truy vấn 3 (Q-P12-03)}
\begin{table}[H]
\centering\small
\caption{Bảng dữ liệu mẫu và kết quả thực thi Truy vấn 3}
\begin{tabularx}{\linewidth}{@{}XXXX@{}}\toprule
\textbf{Ngày đặt} & \textbf{Kênh đặt} & \textbf{Số đơn} & \textbf{Tổng tiền}\\\midrule
2026-10-21 & TaiQuay & 1 & 50000\\
2026-10-20 & TaiQuay & 1 & 52000\\
\bottomrule\end{tabularx}\end{table}
\snapshot{img_06_q03}{Kết quả màn hình hiển thị từ DBMS minh họa thực thi Truy vấn 3}{fig:img_06_q03}
\subsection{Câu truy vấn 4: Số ca và thời lượng trực (Q-P12-04)}
Mô tả: Dùng cho màn hình chấm công và phân ca để thống kê tổng số ca và tổng thời lượng trực của nhân viên; loại bỏ trùng lặp mã ca trước khi cộng tổng.
\sqlinput{sql/Q-P12-04.sql}{Câu lệnh SELECT - Truy vấn 4 (Q-P12-04)}
\begin{table}[H]
\centering\small
\caption{Bảng dữ liệu mẫu và kết quả thực thi Truy vấn 4}
\begin{tabularx}{\linewidth}{@{}XXXX@{}}\toprule
\textbf{Mã NV} & \textbf{Tên nhân viên} & \textbf{Số ca} & \textbf{Số giờ}\\\midrule
TK07 & Demo barista & 2 & 10.00\\
TK01 & Demo admin & 1 & 5.00\\
TK03 & Demo cashier1 & 1 & 5.00\\
\bottomrule\end{tabularx}\end{table}
\snapshot{img_07_q04}{Kết quả màn hình hiển thị từ DBMS minh họa thực thi Truy vấn 4}{fig:img_07_q04}
\subsection{Câu truy vấn 5: Hội viên có số dư vượt mức trung bình (Q-P12-05)}
Mô tả: Dùng cho màn hình chăm sóc khách hàng VIP để tìm hội viên có số dư ví cao hơn mức số dư trung bình của nhóm hội viên từng phát sinh giao dịch ví thành công.
\sqlinput{sql/Q-P12-05.sql}{Câu lệnh SELECT - Truy vấn 5 (Q-P12-05)}
\begin{table}[H]
\centering\small
\caption{Bảng dữ liệu mẫu và kết quả thực thi Truy vấn 5}
\begin{tabularx}{\linewidth}{@{}XXXX@{}}\toprule
\textbf{Mã hội viên} & \textbf{Tên hội viên} & \textbf{Số dư ví} & \textbf{Ngày đăng ký}\\\midrule
TK09 & Demo player3 & 250000 & 2026-10-01\\
\bottomrule\end{tabularx}\end{table}
\snapshot{img_08_q05}{Kết quả màn hình hiển thị từ DBMS minh họa thực thi Truy vấn 5}{fig:img_08_q05}
